\chapter*{chapitre 7}
\markboth{Chapitre 7 }{7 Validation} %pour afficher l'entete
 \addcontentsline{toc}{chapter}{7 Validation}

\textbf{\Huge Validation}

\setcounter{chapter}{7}
\setcounter{section}{0}

\section{Introduction}

Les chapitres précédents ont présenté la réalisation progressive de la plateforme, release après release, en documentant au fil de l'eau les anomalies rencontrées et les corrections apportées. Ce chapitre, qui regroupe le Sprint 11 (tests, validation, corrections et stabilisation) et le Sprint 12 (documentation et préparation de la soutenance), a pour objectif de dresser un bilan global et honnête de l'état de la plateforme au terme du projet : la méthodologie de test suivie, les résultats obtenus, l'évolution de la posture de sécurité mesurée par deux audits successifs, et la synthèse consolidée des limites connues et des perspectives d'amélioration. Conformément à la démarche d'honnêteté académique adoptée tout au long de ce mémoire, ce chapitre ne cherche pas à masquer les points inachevés du projet mais à les documenter avec la même rigueur que les fonctionnalités abouties.


\section{Sprint 11 : Tests, validation, corrections et stabilisation}

\subsection{Introduction}

Ce sprint a pour objectif de valider, de manière transversale à l'ensemble des fonctionnalités développées depuis le Sprint 0, le comportement fonctionnel de la plateforme, et de mener un audit de sécurité formel destiné à identifier et corriger les vulnérabilités les plus critiques avant la mise en exploitation.

\subsection{Méthodologie de test}

La validation fonctionnelle de la plateforme s'est appuyée sur un plan de test manuel structuré, couvrant 84 cas de test répartis sur l'ensemble des modules fonctionnels de la plateforme (authentification et gestion des comptes, gestion des applications, gestion de Kafka, configuration de l'Eventing, facturation, supervision, administration). Ce plan de test constitue le référentiel méthodologique de la campagne de validation ; en toute rigueur, il convient de préciser qu'au moment de la rédaction de ce mémoire, ce document se présente comme une grille de vérification prête à l'exécution, dont les colonnes de statut restent à renseigner de manière exhaustive au fil de son exécution complète — il ne doit donc pas être présenté comme la preuve d'une campagne de test intégralement exécutée et tracée, mais comme la méthodologie de validation retenue pour le projet.

En complément de ce plan de test fonctionnel, la validation technique des fonctionnalités backend non encore exposées graphiquement au moment de leur développement (Kafka et Eventing, avant le Sprint 9) a été réalisée par appels directs à l'API REST à l'aide de l'outil Postman, comme indiqué au chapitre 2.

\subsection{Audit de sécurité et amélioration de la posture de sécurité}

Deux audits de sécurité successifs ont été menés sur la plateforme au cours de ce sprint, permettant de mesurer objectivement la progression de sa posture de sécurité.

Le premier audit a établi un score de sécurité initial de \textbf{4/10}, identifiant notamment les vulnérabilités suivantes, déjà évoquées au fil des chapitres précédents : accès non contrôlé à des ressources appartenant à d'autres tenants (IDOR) sur les journaux de déploiement et sur la suppression de moyens de paiement, secrets applicatifs versionnés en clair, configuration CORS trop permissive combinée à l'acceptation des identifiants de session, absence de limites de ressources par défaut sur les espaces de noms clients, rôles Kubernetes sous-déclarés, transmission du jeton d'authentification dans les paramètres d'URL du flux de journaux, absence de règles d'isolation réseau entre tenants, exécution du déploiement en apparence asynchrone mais en réalité bloquante, et incohérences du flux de journaux en temps réel selon le rôle de l'utilisateur.

Un second audit, mené après correction, a établi un score de \textbf{5{,}2/10}, sur la base du traitement suivant :

\begin{table}[!ht]
\begin{center}
     \begin{tabular}{|p{7cm}|l{5cm}|}
     \hline
     \textbf{Vulnérabilité identifiée} & \textbf{Traitement} \\
     \hline
     IDOR sur les journaux de déploiement & Corrigé \\
     \hline
     IDOR sur la suppression de moyens de paiement Stripe & Corrigé \\
     \hline
     Secrets en clair dans les manifestes versionnés & Corrigé (déplacés en \texttt{Secret} Kubernetes) \\
     \hline
     Configuration CORS permissive avec identifiants de session & Corrigé (liste blanche explicite) \\
     \hline
     Absence de limites de ressources par défaut par tenant & Accepté comme risque produit (modèle \og paiement à l'usage \fg) \\
     \hline
     Rôles Kubernetes (RBAC) sous-déclarés & Corrigé (reconstruits et vérifiés) \\
     \hline
     Jeton d'authentification transmis en paramètre d'URL (SSE) & Corrigé \\
     \hline
     Authentification par mot de passe direct (ROPC) au lieu d'Authorization Code + PKCE & Accepté comme risque produit \\
     \hline
     Absence de règles d'isolation réseau entre tenants & Corrigé (création automatique par tenant) \\
     \hline
     \end{tabular}
     \caption{Synthèse du traitement des vulnérabilités identifiées lors de l'audit de sécurité}
     \label{tab:audit-securite}
     \end{center}
\end{table}

Sept des neuf vulnérabilités critiques ou majeures identifiées ont ainsi été effectivement corrigées, et deux ont fait l'objet d'une décision explicite d'acceptation du risque plutôt que d'une correction, faute de compatibilité immédiate avec le modèle produit visé (facturation à l'usage sans plafond de ressources) ou avec les contraintes de délai du projet (migration du flux d'authentification). Cette distinction entre \og corrigé \fg{} et \og risque assumé \fg{} constitue, de notre point de vue, une démarche d'ingénierie plus rigoureuse qu'une simple liste de correctifs, car elle documente une décision consciente plutôt qu'un oubli.

\subsection{Résultat obtenu}

Au terme de ce sprint, la plateforme a fait l'objet d'une validation fonctionnelle méthodique et d'un durcissement de sécurité mesurable, avec une progression objectivée de 4/10 à 5{,}2/10. La plateforme est jugée stable pour une démonstration et une exploitation pilote, sous réserve des limites explicitement documentées dans la section suivante.

\subsection{Conclusion du Sprint}

Ce sprint a permis de transformer une plateforme fonctionnellement complète en une plateforme dont le niveau de qualité et de sécurité a été objectivement mesuré et amélioré, plutôt que simplement supposé satisfaisant.


\section{Sprint 12 : Documentation et préparation de la soutenance}

\subsection{Introduction}

Ce dernier sprint a pour objectif de consolider la documentation technique de la plateforme et de préparer sa présentation devant le jury.

\subsection{Réalisation}

L'ensemble des décisions techniques, anomalies identifiées et correctifs apportés tout au long du projet ont été consignés de manière détaillée dans une documentation technique versionnée avec le code (guide technique complet, fiches de correctifs individuelles, rapports d'audit), constituant la source principale de ce mémoire. Un plan de sauvegarde et de restauration de la base de données PostgreSQL et de l'index Elasticsearch a également été formalisé, avec une politique de rétention de sept sauvegardes quotidiennes et quatre sauvegardes hebdomadaires.

\subsection{Résultat obtenu}

La plateforme dispose, au terme de ce sprint, d'une documentation technique exhaustive et d'un mémoire retraçant fidèlement sa réalisation, ses choix d'architecture et ses limites, prêts à être présentés devant le jury.

\subsection{Conclusion du Sprint}

Ce sprint clôt la conduite du projet selon la planification Scrum définie au premier chapitre.


\section{Synthèse des limites et perspectives d'amélioration}

Cette section consolide, en un point unique, l'ensemble des limites identifiées au fil des chapitres de réalisation, ainsi que les fonctionnalités explicitement non implémentées à ce stade du projet.

\begin{table}[!ht]
\begin{center}
     \begin{tabular}{|p{6.5cm}|p{7cm}|}
     \hline
     \textbf{Limite ou fonctionnalité absente} & \textbf{Statut / perspective} \\
     \hline
     Infrastructure du cluster (nœuds, CNI, MetalLB, Kourier) et cluster Kafka Strimzi non versionnés en infrastructure-as-code & Provisionnés manuellement (Sprints 0 et 2) ; perspective : formalisation via Helm/Kustomize/GitOps \\
     \hline
     Absence de limites de ressources (\texttt{LimitRange}) par défaut par tenant & Risque produit assumé (modèle à l'usage) \\
     \hline
     Flux d'authentification ROPC au lieu d'Authorization Code + PKCE & Risque produit assumé \\
     \hline
     Absence d'isolation forte (SASL/mTLS/ACL) entre tenants sur le cluster Kafka partagé & Non implémenté ; perspective : \texttt{KafkaUser} par tenant avec ACL dédiées \\
     \hline
     Jeton et jeton de rafraîchissement stockés dans le stockage local du navigateur & Non corrigé ; perspective : migration vers un stockage plus sûr (cookie \textit{httpOnly}) \\
     \hline
     Collecte Prometheus des métriques du backend potentiellement toujours inopérante (étiquette de sélection incohérente) & Anomalie identifiée, correction non confirmée en production \\
     \hline
     Console d'administration exposée sur un port réseau fixe sans protection réseau dédiée & Non corrigé ; perspective : restriction d'accès réseau ou VPN dédié \\
     \hline
     Absence d'analyse de vulnérabilités des images de conteneurs dans les pipelines CI/CD & Non implémenté ; perspective : intégration d'un scanner d'images (ex. Trivy) \\
     \hline
     Conteneurs applicatifs exécutés avec l'utilisateur root par défaut & Non corrigé ; perspective : imposition d'utilisateurs non privilégiés \\
     \hline
     Pipelines CI/CD hétérogènes (portail, console, microservices moins aboutis que le backend) & Partiellement traité au Sprint 10 ; perspective : socle Jenkins commun \\
     \hline
     Absence de tests automatisés (unitaires/intégration) dans les pipelines CI/CD & Non implémenté ; perspective : intégration d'une étape de tests et de couverture de code \\
     \hline
     Duplication de code entre le portail client et la console d'administration & Non refactorisé ; perspective : extraction d'un paquet partagé ou passage à un monorepo \\
     \hline
     Mise en cache et découpage par domaine personnalisé (\textit{traffic splitting}, \textit{domain mapping}), non retrouvés dans le code audité & Non implémenté ; perspective d'évolution \\
     \hline
     \end{tabular}
     \caption{Synthèse des limites connues et des perspectives d'amélioration}
     \label{tab:limites}
     \end{center}
\end{table}

Cette synthèse illustre une caractéristique assumée de la démarche suivie tout au long du projet : privilégier, à ressources et délai constants, la robustesse fonctionnelle du parcours principal (déploiement, supervision, facturation) plutôt qu'une couverture exhaustive de toutes les dimensions non fonctionnelles souhaitables. Les limites identifiées ne remettent pas en cause la validité de la démonstration technique réalisée, mais constituent une feuille de route réaliste pour une éventuelle industrialisation ultérieure de la plateforme.


\section{Conclusion}

Ce chapitre a présenté la validation de la plateforme au travers d'une campagne de tests fonctionnels méthodique et d'un audit de sécurité en deux temps, ayant permis de mesurer objectivement une amélioration de la posture de sécurité du système. Il a également dressé la synthèse honnête des limites connues du projet, condition nécessaire, à notre sens, à la crédibilité scientifique et professionnelle de ce mémoire devant un jury de cycle ingénieur. La conclusion générale qui suit revient sur les apports globaux du projet au regard de la problématique posée au premier chapitre.
